\documentstyle[%


righttag%


,11pt%


,amssymb%


,russian%               remove in Eng.


,izvru%                 Set izven% in Eng.


]{amsart}





%





\numberwithin{equation}{section}


\newcommand{\e}{\varepsilon}


\newcommand{\g}{\gamma}


\newcommand{\f}{\varphi}


\renewcommand{\l}{\lambda}


\renewcommand{\o}{\omega}


\renewcommand{\d}{\partial}


\renewcommand{\aa}{\alpha}


\renewcommand{\t}{\widetilde}


\renewcommand{\v}{\upsilon}


\newcommand{\tts}[1]{}


\renewcommand{\baselinestretch}{1.05}





\theoremstyle{plain}


\theorembodyfont{\it}


\newtheorem{thmnonum}{\hskip\parindent Теорема}


\renewcommand{\thethmnonum}{}





%\hoffset=-15mm%                  For Eng.


\setcounter{page}{3}


\god{2002}


\nomer{3~(478)} %                        Remove in Eng.


%\nomer{Vol.\,46, No.\,3}%               Set in Eng.


%\str{\pageref{begin}--\pageref{end}}%  Set in Eng.


\udk{517.9}





\author{\it В.Н.\,Бобочко}


% NB:   ^^ remove in Eng.


\title{Дифференциальная точка поворота\\


в теории сингулярных возмущений. I}





\date{}


%\pagestyle{myheadings}%         Set in Eng.


\pagestyle{plain} %              Remove in Eng.


\begin{document}


%\thispagestyle{plain}%          Set in Eng.


\maketitle


\label{begin}








\section*{Введение}





Теория сингулярно возмущенных дифференциальных 


уравнений (СВДУ) с точками поворота


начала свое бурное развитие с исследований Р.\,Лангера [1]--[4].


В настоящее время основные результаты по построению равномерных 


асимптотик для уравнения Лиувилля и систем


СВДУ типа Лиувилля с точкой поворота уже получены [1]--[14]. 


Характерной чертой большинства этих


работ было то, что вырожденные уравнения были алгебраическими,


т.\,е. точка поворота была алгебраической.





Однако теория и практика требуют исследований СВДУ, для которых


вырожденные уравнения являются дифференциальными и точка поворота 


находится в виде множителя при старшей производной вырожденного уравнения.





Классическими примерами таких уравнений являются уравнения


\begin{gather}


\bold{L}_{\e} y (x, \e) \equiv \e^3 y''' (x, \e) +x \wt{a}


(x) y' (x, \e) + b(x) y(x, \e) = h (x), \tag{0.1} \\


%


\e y'' (x, \e) + x \wt{a} (x) y' (x, \e) + b(x) y (x, \e) =h(x) \tag{0.2}


\end{gather}


и уравнение Орра-Зоммерфельда


\begin{equation}


\e y^{(4)} (x, \e) + x \wt{a} (x) y'' (x, \e) + b(x) y' (x, \e)


+ c(x) y(x, \e) = h(x).\tag{0.3}


\end{equation}


Общей чертой всех этих трех классических уравнений является то, что


все корни соответствующих характеристических уравнений совпадают в


точке $x=0$, т.\,е. в точке поворота. В то же время каждое из этих


уравнений имеет существенные особенности, присущие только ему.


Исследование и построение равномерных асимптотик решений этих


уравнений является сложной задачей. Проблема построения


равномерных асимптотик решений таких уравнений насчитывает


несколько десятилетий (см. [4]; [13], с.\,225--230). Еще в 60-е


годы С.А.\,Ломов перед своими учениками поставил задачу по


исследованию уравнения (0.2) (задача Гордеева). Однако


фундаментальных результатов в этом направлении до настоящего


времени так и не получено. Здесь уместно напомнить, что методом


согласования в ([14], с.\,52--63) построена асимптотика


решения уравнения (0.2) на отрезке $[0; 1]$.





Для уравнений (0.1) и (0.3) в [2]--[4] тоже были


получены определенные результаты. Однако они носили частный


характер, и дальнейших обобщений этих результатов не было получено.





Цель данной работы будет состоять в обобщении метода,


разработанного в [7]--[12] для исследования скалярных и векторных


уравнений с алгебраической точкой поворота на уравнения с


дифференциальной точкой поворота. Первым шагом в этом направлении


и будет исследование сингулярно возмущенного дифференциального


уравнения (0.1).





Уравнение (0.1) иcследуется для случая, когда выполняются





{\it Условие} $1^0$. $\quad$ $a(x), b(x), h(x) \in


\bold{C}^{\infty} [I]$, $I=[0; 1]$;





{\it Условие} $2^0$. $\quad$ $a(x) = x \wt{a} (x)$, причем


$\wt{a} (x) >0$, $b(x) < 0$ при $x \in I$.





\section{Структура решения вырожденного уравнения}





{\it Условие}~$2^0$ обеспечивает существование достаточно гладкого


решения уравнения 


\begin{equation}


\bold{L}_0 \o (x) \equiv x \wt{a} (x) \o' (x) + b(x) \o (x) =h(x).


\end{equation}


Поскольку решение этого уравнения будет играть важную роль в


построении асимптотики решения СВДУ (0.1), то исследуем структуру


этого решения, особенно в точке поворота.





Разлагая в ряд Маклорена, получим равенства


\begin{equation}


- \frac{b(x)}{x \wt{a} (x)} = \frac{\rho}{x} + \aa (x), \quad


\frac{h(x)}{\wt{a} (x)} = q + x \g (x),


\end{equation}


где $\aa(x)$ и $\g (x)$ --- аналитические функции в окрестности точки


$x=0$, причем $\g (0) \neq 0$ и $\rho = \frac{-b (0)}{\wt{a}(0)}> 0$.


С учетом равенств (1.2) решение вырожденного уравнения (1.1)


запишем в виде


\begin{equation}


\o (x) = x^{\rho} \beta (x) \bigg\{ C + \int _1^x [ q + x \g (x)


] x^{-\rho -1} \beta^{-1} (x) d x \bigg\}, 


\end{equation}


где $C$ ---


произвольная постоянная, а $\beta (x) = \exp \{ \int \limits_1^x


\aa (x) d x \}$.





Легко проверить, что при $\rho > 0$ решение (1.3) является


непрерывной функцией при всех $x \in [0; 1]$, причем


$\o (0) =-\rho^{-1} q < \infty$.





Можно было бы выписать и начальные значения $\o^{(s)} (0)$,


$s=\overline{1, 3}$, однако явной необходимости в этом нет.


Заметим только, что эти значения уже зависят от принимаемых


значений параметра $\rho$. Для обеспечения построения равномерно


пригодной асимптотики решения СВДУ (0.1) на всем отрезке $I =[0,1]$


с любой степенью точности относительно малого параметра


необходимо будет потребовать, чтобы $ \rho \in N $. Если же $\rho>0$


достаточно большое, не является натуральным числом, то легко


просчитать с какой степенью точности можно построить асимптотику


решения исследуемого СВДУ (0.1).





\section{Расширение возмущенного уравнения}





Особая точка $\e = 0$ порождает в решении СВДУ (0.1) некоторые существенно


особые функции (СОФ). Основная трудность построения асимптотики


решения СВДУ (0.1) состоит в определении этих СОФ. С этой целью


постараемся сохранить и применить общий метод исследования СВДУ с


алгебраической точкой поворота, разработанный в [7]--[12].





Hаряду с независимой переменной $x \in I$ введем 


новую переменную $t$ по формуле


\begin{equation}


t = \e^{-p} \f (x) \equiv \Phi (x, \e), 


\end{equation}


где показатель $p$ и регуляризующая функция $\f (x)$


подлежaт определению.





Учитывая (2.1), согласно


методу регуляризации будем исследовать расширенную функцию


$\wt{y} (x, t, \e)$, для которой должно выполняться тождество


\begin{equation}


\wt{y} (x, t, \e)


\big|_{t= \e^{-p} \f (x)} \equiv y(x,\e). 


\end{equation}


Дифференцируя тождество (2.2), определим полные производные до


третьего порядка и подставим значения этих производных в СВДУ


(0.1). Тогда для определения расширенной функции $\wt{y} (x, t,


\e)$ получим расширенное уравнение


\begin{equation}


\wt{\bold{L}}_{\e} \wt{y} (x, t, \e) = h(x), 


\end{equation}


в котором расширенный оператор $\wt{\bold{L}}_{\e}$ имеет вид


\begin{multline}


\wt{\bold{L}}_{\e} \equiv \e^{3-3p} {\f'}^3 (x)


\frac{\d^3}{\d t^3} + \e^{3-2p} \bigg[3 {\f'}^2 (x) \frac{\d}{\d x} + 3


\f' (x) \f'' (x) \bigg] \frac{\d^2}{\d t^2} + \e^{3-p} \bigg[ 3 \f' (x)


\frac{\d^2}{\d x^2 } + \\


%


+ 3 \f'' (x) \frac{\d}{\d x} + \f''' (x) \bigg] \frac


{\d}{\d t} + a(x) \bigg[ \e^{-p} \f' (x) \frac{\d}{\d t} + \frac{\d}{\d


x} \bigg] + \e^3 \frac{\d^3}{\d x^3} + b(x).


\end{multline}





Для дальнейших исследований, т.\,е. для описания пространства


безрезонансных решений (ПБР) необходимо выделить в расширенном


операторе (2.4) некоторый модельный оператор. Для алгебраических


точек поворота модельными операторами были классические операторы


Эйри, т.\,е. операторы $\frac{\d^2}{\d t^2} \pm t$, свойства которых


описаны в [5], [13].





Если для исследования решения СВДУ (0.1) применить классический


оператор Эйри, то в процессе построения коэффициентов асимптотики


решения получим дифференциальное уравнение с иррегулярной особой


точкой $x=0$. А это означает, что построенная таким образом


асимптотика решения будет неограниченно возрастать, когда $x \to+0$,


т.\,е. неограниченно возрастает в окрестности точки поворота.





Для построения асимптотики решения однородного уравнения (0.1) в


[2]--[3] использовано модельное уравнение


\begin{equation}


\frac{\d^2 \upsilon}{\d z^2} + \l^2 z \frac{\d \upsilon}{\d z} +


3 \l \mu \upsilon =0, 


\end{equation}


где $| \l | \to + \infty $, $\mu$ --- некоторое число. Такого


же типа модельный оператор был применен в ([14], с.\,55)


при исследовании краевой задачи для уравнения (0.2).





Проведенные автором исследования показали, что для построения


равномерной асимптотики решения СВДУ (0.1) есть смысл использовать


модельный оператор 


\begin{equation}


\wt{\bold{T}} \equiv \frac{\d^3}{\d t^3} + t \frac{\d}{\d t} \equiv


\bold{T}


\frac{\partial}{\partial t}, 


\end{equation}


который будем называть { \it дифференциальным оператором


Эйри--Дородницына}.





\begin{note}


Вместо модельного оператора (2.6) можно было взять оператор


$\bold{T} \equiv \frac{\d^3}{\d t^3} - t \frac{\d}{\d t}$,


т.\,е. дифференциальный оператор Эйри--Лангера.


\end{note}





Оператор (2.6) выделим в расширенном операторе (2.4) следующим


образом:


$$


\e^{3-3p} {\f'}^3 (x) \frac{\d^3}{\d t^3} + \e^{-p} \f'(x) a(x)


\frac{\d}{\d t} \equiv \e^{3-3p} {\f'}^3 (x) \bigg[ \frac{\d^3}{\d t^3}


+ \frac{a(x)}{{\f'}^2 (x) \e^{3-2p}} \frac{\d}{\d t} \bigg] \equiv


\e^{3-3p} {\f'}^3 (x) \bold{T},


$$


где


\begin{equation}


\frac{a(x)}{{\f'}^2 (x) \e^{3-2p}}


\equiv t \big|_{t=\e^{-p} \f(x)} \equiv \e^{-p} \f (x). 


\end{equation}





Используя (2.7), определим показатель $p=1$ и получим следующее


дифференциальное уравнение для определения регуляризующей функции:


\begin{equation}


[{\f'} (x)]^2 \, \f (x) = a(x) \equiv x \wt{a} (x). 


\end{equation}


Решением уравнения (2.8) при начальном условии $\f (0) = 0$ будет


функция


$$


\f (x) = \biggl(\frac{3}{2} \int_0^x \sqrt{x \wt{a}


(x)} d x {\biggr)}^{2/3}. 


$$





В исследуемом случае $\f (x) \geq 0$ при $x \in I$, т.\,е. $t \geq0$.


Это условие в будущем обеспечит ограниченность СОФ


$W_k^{(s)} (t)$, $k = 1,2$, $s = 0, 1$, (см. (3.2)), содержащихся в


решении уравнения (0.1). Точку поворота $x=0$, порождающую


ограниченность вышеупомянутых СОФ, будем называть {\it стабильной}


точкой поворота.





Если же $\wt{a} (0) < 0$, $x \in I$, то СОФ $ U_2 (t)$ и все


его производные неограниченно возрастают при $t \rightarrow +


\infty$. Проведенные исследования показали, что в этом случае


асимптотика решения и особенно оценка остаточного члена


асимптотики решения в определенной степени отличаются от


исследуемого случая, т.\,е. когда $ \wt{a} (0)> 0$.


В этом случае точку поворота $x=0$ будем называть {\it


нестабильной} точкой поворота. Случай алгебраической нестабильной


точки поворота тоже исследован автором (см. [10] и другие работы


автора).





\begin{note}


Случаи, когда коэффициенты уравнения (0.1)


удовлетворяют условиям, отличным от {\it условия} $2^0$,


будут предметом последующих исследований.


\end{note}





С учетом тождества (2.6) и определенном показателе $p = 1$


расширенный оператор (2.4) запишем в виде


\begin{equation}


\wt{\bold{L}}_{\e} \equiv {\f'}^3 (x) \wt{\bold{T}} +


\bold{L}_0 + \e\, \bold{d} \frac{\d^2}{\d t^2 }


+ \e^2\, \bold{m}\, \frac{\d}{\d t} + \e^3 \frac{\d^3}{\d x^3},


\end{equation}


где


\begin{gather*}


\bold{L}_0 \equiv a(x) \frac{\d}{\d x} + b(x), \\


\bold{d} \equiv 3 {\f'}^2 (x) \frac{\d}{\d x}


+ 3 \f'(x) \f''(x), \quad \bold{m} \equiv 3 \f' (x)


\frac{\d^2}{\d x^2} + 3 \f''(x) \frac{\d}{\d x} + \f'''(x).


\end{gather*}





\section{Пространства безрезонансных решений}





С расширением и введением модельного оператора $\wt{\bold{T}}$ был


в значительной степени определен характер асимптотики решения


исследуемого СВДУ (0.1). Для полной ясности в этом вопросе


необходимо еще описать и исследовать ПБР, в которых будем строить 


асимптотику решения уравнения (0.1).





Сначала исследуем модельное уравнение


\begin{equation}


\wt{\bold{T}} W(t) = 0. 


\end{equation}


Двумя линейно независимыми решениями этого уравнения будут функции


\begin{equation}


W_k (t) \equiv \int_{\infty}^t U_k (\tau) d \tau, \quad k=1, 2,


\end{equation}


где $U_k (t)$ --- функции Эйри--Дородницына, т.\,е. два


линейно независимыx решения модельного уравнения $U''(t) + t U(t)= 0$.


Эти решения удовлетворяют начальным условиям


\begin{gather}


W_1 (0) = - 3^{{-1}/{3}} \Gamma (2/3)


, \quad {W'}_1 (0) = 1,\quad {W''}_1 (0) = 0, \quad W_2 (0) = - 3^{{1}/{3}}


\Gamma (4/3), \notag \\


%


{W'}_2 (0) = 0, \quad {W''}_2 (0) = 1.


\end{gather}


Третьим линейно независимым решением уравнения (3.1) будет


функция $U_3 (t)=1$.





Поскольку $W_k^{(s)} (t) = U_k^{(s-1)} (t)$, то асимптотические


свойства производных функций (3.2) известны [5]. В этом и


состоит одно из преимуществ использования модельного уравнения


(3.1) по сравнению с модельным уравнением (2.5). Для получения


асимптотических формул для самих функций (3.2) при $ t \rightarrow


+ \infty $ поступим следующим образом. Подставим в (3.2)


соответствующие формулы для функций Эйри и проинтегририруем по


частям. В результате получим следующие асимптотические равенства


при $t \rightarrow + \infty$:


\begin{align*}


W_1 (t) &= A \,t^{{-3}/{4}} \biggl\{ \sin \aa [1 + O \bigl(t ^


{-3/2} \bigr) ] - \frac{41}{48} t^{{-3}/{2}} \cos \aa [1 +


O \bigl(t^{{-3}/{2}} \bigr) ] \biggr\},\\


%


W_2 (t) &= B\, t^{{-3}/{4}} \biggl\{ \cos \beta [1 + O \bigl(t^{-3/2}


\bigr) ] + \frac{41}{48} t^{{-3}/{2}} \sin \beta [1 +


O \bigl(t^{{-3}/{2}} \bigr) ] \biggr\},


\end{align*}


где $A = \pi^{{-1}/{2}}\, 3^{{1}/{6}}\, \Gamma \left(\frac{2}{3} \right)$,


$B = \pi^{{-1}/{2}}\, 3^{{-1}/{6}}\, \Gamma \left(\frac{1}{3} \right)$,


$\aa = \frac{2}{3}\, t^{{3}/{2}}-\frac{\pi}{12}$,


$\beta = \frac{2}{3}\, t^{{3}/{2}}+\frac{\pi}{12}$.





Представим решения (3.2) в окрестности точки $t = 0$ в виде


степенных рядов по переменной $t$ и потребуем, чтобы они


удовлетворяли начальным условиям (3.3). Тогда получим


представления


\begin{align*}


W_1 (t) & \equiv -3^{{-1}/{3}}\, \Gamma \left(\frac{2}{3} \right) + t -


\frac{t^4}{4 \cdot 3 \cdot 2} + \frac{4}{7\cdot 6 \cdot 5 }


\cdot \frac{1}{4 \cdot 3 \cdot 2 } t^7 + O (t^{10}), \\


%


W_2 (x,t) & \equiv -3^{{-1}/{3}}\, \Gamma \left(\frac{4}{3}\right)


+ \frac{1}{2}\, \biggl[ t^2 - \frac{2t^5}{5 \cdot 4 \cdot 3}


+\frac{5}{8\cdot 7 \cdot 6 } \cdot \frac{2}{5 \cdot 4 \cdot 3 } t


^8 +O (t^{11}) \biggr].


\end{align*}





Изучив свойства СОФ (3.2), перейдем к описанию ПБР.





Используя (3.2), введем множества (подпространства) функций


\begin{equation}


Y_{rk} = \{ V_{rk} (x) W_k (t) + Q_{rk} (x) U_k (t) + M_{rk} (x) U'_k (t) \},


\ \  k=1, 2, \quad X_r = \{ \o_r (x) \}, 


\end{equation}


где $V_{rk} (x)$, $Q_{rk} (x)$, $M_{rk} (x)$, $\o_{r} (x) \in


\bold{C}^{\infty} [I]$.


Из подпространств (3.4) составим новое пространство


\begin{equation}


Y_r = \bigoplus \limits_{k=1}^2 Y_{rk} \bigoplus X_r, 


\end{equation}


которое согласно принятой терминологии (см. [6], с.\,42), будем называть ПБР.


Элемент ПБР (3.5) имеет вид


\begin{equation}


y_r (x, t) = \sum \limits_{k=1}^2 \bigl[ V_{rk} (x) W_k (t) +


Q_{rk} (x) U_k (t) + M_{rk} (x) U'_k (t) \bigr] + \o_r (x)


\equiv \sum \limits_{k=1}^2 y_{rk}(x,t) + \omega_r (x). 


\end{equation}





Согласно разработанному общему методу (см. [6]--[12]) необходимо


изучить действие расширенного оператора $\bold{\wt{L}}_{\e}$


(см. (2.9)) на элемент ПБР (3.5), т.е. на элемент (3.6).





Если ПБР (3.5) описано правильно, т.\,е. если в этом пространстве


правильно выделены и сохранены как единые целые все СОФ,


содержащиеся в решении СВДУ (0.1), то в ПБР расширенное уравнение


уже должно быть регулярно зависящим от малого параметра $\e > 0$.





Используя модельный дифференциальный


оператор (2.6), легко получить тождества


\begin{gather}


W'''_k (t) \equiv - t W'_k (t)


\big|_{t= \e^{-1} \f(x)}\equiv


-\e^{-1} \f(x) W'_k (t) \equiv -\e^{-1} \f(x) U_k (t), \notag \\


%


W_k ^{(4)} (t) \equiv - W'_k (t) - t W''_k (t)


\big|_{t= \e^{-1} \f(x)}\equiv -U_k (t) -


\e^{-1} \f (x) U'_k (t), \\


%


W_k ^{(5)} (t) \equiv -2 W''_k (t) + t^2 W'_k (t)


\big|_{t= \e^{-1} \f(x)}\equiv


-2 U'_k (t) + \e^{-2} \f^2 (x) U_k (t);\notag \\


%


\wt{\bold{T}} W_k (t) \equiv 0, \quad \bold{\wt{T}} W'_k (t) \equiv -


W'_k (t) \equiv - U_k (t), \quad \bold{\wt{T}}


W''_k (t) \equiv -2{W''}_k (t) \equiv -2U'_k (t).


\end{gather}


Заметим, что частичное сужение, проведенное в тождествах (3.7), в


будущем обеспечит отсутствие резонансных множителей $t$ в


построенном асимптотическом решении СВДУ (0.1), т.\,е. будет получен


регуляризованный асимптотический ряд решения (см. [6], с.\,39).





С учетом тождеств (3.7) и (3.8) получим результат действия


расширенного оператора $\bold{\wt{L}}_{\e}$ (см. (2.9)) на


элемент $y_r (x, t) \in Y_r $, представимый в виде тождества


\begin{equation}


\wt{\bold{L}}_{\e} y_r (x, t) \equiv \sum \limits_{s=0} ^3


\e^s \bold{R}_s y_r (x, t)\equiv \sum \limits_{s=0} ^3 \e^s \sum


\limits_{k=1} ^2 \bold{R}_{sk} y_{rk} (x, t)


+ \wt{\bold{L}}_{\e} \omega_s (x), 


\end{equation}


где 


\begin{gather}


\bold{R}_0 y_{r} (x, t) \equiv \sum \limits_{k=1}^2 \biggl\{ \bold{L}_0


V_{rk} (x) W_k (t) + \bold{D}_1 Q_{rk} (x)U_k (t)


+ \bold{D}_2 M_{rk} (x) U'_k (t) \biggr\}


+\bold{L}_0 \o_r (x), \\


%


\bold{R}_1 y_r (x, t) \equiv \sum \limits_{k=1}^2 \biggl\{ -


\bigl(\bold{d} + \bold{m}\, \f(x) \bigr) M_{rk} (x)\, U_k (t)


+ \bold{d}\, V_{rk} (x) \,U_k\,' (t) \biggr\} \equiv \sum


\limits_{k=1}^2 \bold{R}_{1k} y_{rk} (x, t), \\


%


\bold{R}_2 y_r (x, t)  \equiv \sum \limits_{k=1}^2 \biggl\{


\bold{m} \, V_{rk} (x) U_k (t) + \bold{m} \, Q_{rk} (x) U_k\,'


(t) \biggr\}, \quad \bold{R}_3 y_r (x, t) \equiv \frac{\d^3 y_r


(x, t)}{\d x^3},


\end{gather}


\begin{align}


\bold{D}_1 & \equiv


\bold{L}_0 - [{\f'} (x)]^3 - \f(x)\, \bold{d} \equiv -2a(x)


\frac{\d}{\d x} + b_1 (x), \\


%


\bold{D}_2 & \equiv


\bold{L}_0 - 2 [{\f'}(x)]^3 - \f(x)\, \bold{d} \equiv -2a(x)


\frac{\d}{\d x} + b_2 (x),


\end{align}


$$


b_1 (x) \equiv b(x) - {\f'}^3 (x) - 3\f(x) \f'(x) \f''(x), \quad


b_2 (x) \equiv b(x) -2 {\f'}^3 (x) - 3\f(x) \f'(x) \f''(x).


$$





Из полученных тождеств (3.9)--(3.14) можно сделать следующие выводы.





1. Пространства безрезонансных решений (3.5) инвариантны


относительно операторов $\bold{R}_s$, $s=\overline{0, 3}$, а


следовательно, и относительно расширенного оператора


$\bold{\wt{L}}_{\e}$, представимого в виде тождества (3.9).





2. Оператор $\bold{R}_0$ является главным оператором


расширенного оператора $\bold{\wt{L}}_{\e}$ в ПБР (3.5).





3. Каждое из подпространств $Y_{r1}$, $Y_{r2}$ и $X_r$ инвариантно


относительно расширенного оператора $\bold{\wt{L}}_{\e}$.





4. Расширенное уравнение (2.3) является регулярно возмущенным относительно


малого параметра $\e >0$ в ПБР $Y_r$.





Следовательно, здесь проведена {\it регуляризация} СВДУ (0.1). Это


один из основных пунктов метода регуляризации существенно особых


функций (метод РСОФ).








\section{Формализм построения ряда решений расширенного уравнения}





На основании сделанных выводов асимптотику решения расширенного


уравнения (2.3) строим в виде ряда


\begin{equation}


y (x, t, \e) = \sum \limits_{r=0}^{+\infty} \e^r y_r (x, t) \equiv


\sum \limits_{r=0}^{+\infty} \e^r \bigg[ \sum \limits_{k=1}^{2}


y_{rk} (x, t) + \omega_r (x) \bigg], \quad y_r (x,t) \in Y_r.


\end{equation}





Для определения коэффициентов ряда (4.1) получим следующую


рекуррентную систему уравнений:


\begin{gather}


\bold{R}_0 y_0 (x, t) = h(x), \\


\bold{R}_0 y_1 (x, t) = - \bold{R}_1 y_0 (x, t), \\


\bold{R}_0 y_2 (x, t) = - \bold{R}_1 y_{1} (x, t) -


\bold{R}_2 y_{0} (x, t),\notag \\


\bold{R}_0 y_r (x, t) = - \bold{R}_1 y_{r-1} (x, t) -


\bold{R}_2 y_{r-2} (x, t) - \bold{R}_3 y_{r-3} (x, t), \quad r \geq 3. 


\end{gather}








\section{Построение общего решения расширенного уравнения}





Покажем, что





1) серия рекуррентных уравнений (4.2)--(4.4)


асимптотически корректна в


ПБР, т.\,е. каждое из этих уравнений имеет решение в ПБР


(3.5) при последовательном решении всех предыдущих уравнений;





2) решение каждого уравнения (4.2)--(4.4) является


достаточно гладким на всем отрезке ${I=[0; 1]}$, т.\,е. и в точке


поворота $x=0$;





3) каждое решение $y_r (x,t)$ содержит по три произвольные постоянные.





Приступим к исследованию решения первого итерационного уравнения


(4.2) в ПБР $Y_0$.


В силу линейной независимости СОФ $W_k (t)$ и их производных


в ПБР $Y_0$ уравнение (4.2) распадается на следующие


дифференциальные уравнения ($r=0)$:


\begin{gather}


\bold{L}_0 V_{rk} (x) = 0, \quad k = 1,2, \quad \bold{L}_0


\o_r (x) = h(x), \\


\bold{D}_1 Q_{rk} (x) = 0, \quad \bold{D}_2 M_{rk} (x) =0,


\quad k = 1,2.


\end{gather}


Из п.\,1 (см. (1.1)--(1.3)) следует, что существуют


достаточно гладкие решения дифференциальных уравнений $ (5.1)$ на


всем отрезке $[0.1]$, представимые при $r = 0$ в виде


\begin{gather}


V_{rk} (x) = C_{rk} x^{\rho} \beta (x)+ \wt{V}_{rk} (x) \equiv


C_{rk} x^{\rho} \exp \biggl\{ \int _1^x \aa(x) d x


\biggr\}+ \wt{V}_{rk} (x), \\


%


\o_{r} (x) = x^{\rho} \beta (x) \biggl[ C_{r3} + \int _1^x


\biggl[ q + x \gamma (x) \biggr] x^{\rho -1} \beta^{-1} (x) d x


\biggr] \equiv C_{r3} x^{\rho} b (x) + \wt{ \o}_r (x) +


\wt{\o}_r^* (x),


\end{gather}


где $C_{0s}$, $s=\overline{1,3}$, --- произвольные 


постоянные, а $\wt{\o}_r^* (x)$


--- частные решения, получены за счет неоднородного дифференциального


уравнения $(5.1)$, причем они будут


тождественными нулями в случае, когда $h (x) \equiv 0$. Слагаемые


$ \wt{V}_{rk} (x)$ и $\wt{\o}_r (x)$ возникают за счет


интегрирования неоднородных дифференциальных


уравнений с правыми частями, содержащими произвольные постоянные


от предыдущих решений. Поэтому эти функции в отличиe от $\wt{\o}_r^* (x)$


участвуют в построении линейно независимых решений СВДУ (0.1).





Полученные решения (5.3) и (5.4) уже содержат три произвольные постоянные


$C_{0s}$, $s=\overline{1,3}$.


Поэтому для построения общего решения расширенного уравнения (2.3),


а следовательно, и СВДУ (0.1), достаточно


строить только частные решения дифференциальных


уравнений $(5.2)$. Операторы $\bold{D}_i$, $i=1, 2$, имеют особенности при


производных. Поэтому гладкими решениями однородных


дифференциальных уравнений $(5.2)$ при условиях $| Q_{0k} (0)| <\infty$


и $|M_{0k} (0)| < \infty$ будут тождественные нули.





Линейная независимость


функции $W_k (t)$, $k=1, 2$, в дальнейшем обеспечивает линейную


независимость частных решений $y_k (x, \Phi(x,\e), \e)$, $k=1,2$,


исследуемого уравнения (0.1).





Таким образом, достаточно гладким решением на всем отрезке $[0;


1]$ неоднородного уравнения (4.2) в ПБР $Y_r$ будет функция


\begin{equation}


y_0 (x, t) = \biggl[ \sum \limits_{k=1}^2 C_{0k} W_k (t) + C_{03}


\biggr] x^\rho \beta(x) + \wt{\o}_r^* (x) \equiv \sum


\limits_{s=1}^3 Y_{0s} (x, t) + \wt{\o}_r^* (x).


\end{equation}





Переходим к решению итерационного уравнения (4.3). Сначала


вычислим правую часть этого уравнения. Используя тождество


(3.11), имеем $- \bold{R}_1 y_0 (x, t) = - \sum \limits_{k=1}^2


\bold{d} \,V_{0k} (x) U'_k (t)$. По аналогии с предыдущим для


определения коэффициентов решения $y_1 (x, t)$ получим однородные


уравнения (5.1) при $r=1$ и дифференциальные уравнения


\begin{equation}


\bold{D}_1 Q_{1k} (x) =0, \quad \bold{D}_2 M_{1k} (x) = -


\bold{d} \, V_{0k} (x),\quad k=1,2. \tag{$5.2^1$}


\end{equation}


Снова получим $Q_{1k} (x) \equiv 0$. Второе уравнение $(5.2^1)$


уже неоднородное. Поэтому существует не равное тождественно нулю


достаточно гладкое на $[0; 1]$ решение этого уравнения.


Следовательно, общим решением неоднородного итерационого уравнения


(4.3) будет


$$


y_1 (x, t) = \biggl[ \sum \limits_{k=1}^2 C_{1k} W_k (t) + C_{13}


\biggr] x^{\rho} \beta(x) + \sum \limits_{k=1}^2 M_{1k} (x)


U_k\, ' (t) \equiv \sum \limits_{k=1}^3 Y_{1k} (x, t).


$$





Продолжим далее решать итерационные уравнения (4.4). В процесе


постепенного решения этих уравнений будут определены функции


$$


y_r (x, t) = \sum _{k = 1}^2 Y_{rk} (x, t) + Y_{r3} (x) +\wt{\omega}_r^* (x),


$$


где $\wt{\o}_r^* (x)$ --- частные решения


дифференциальных уравнений, полученные от интегрирования


неоднородного СВДУ (0.1), т.\,е. $\wt{\o}_r^* (x)\equiv 0 $ при


$h(x)\equiv 0$.





\begin{note}


Условиe $2^0$, а именно, противоположность


знаков функций $\wt{a}(x)>0$ и $b(x)~<~0$ при $x \in I$,


позволило получить гладкое решение уравнения (1.1), причем


зависящее от произвольной постоянной. В дифференциальных


операторах $\bold{D_1}$ и $\bold{D_2}$ (см. (3.13)--(3.14)) в


общем случае это условие уже не сохраняется. Если бы


это свойство сохранялось и для операторов $\bold{D_1}$ и


$\bold{D_2}$, то для построения равномерной асимптотики решения


СВДУ (0.1) достаточно взять только частные решения


дифференциальных уравнений (5.2).


\end{note}





Проведенные исследования


показали, что три решения однородного расширенного уравнения


(2.3) можно построить в виде формальных


рядов


\begin{equation}


\negthickspace


Y_k (x, t, \e) = \sum \limits_{r=0}^{\infty} \e^r \biggl[ V_{rk}


(x) W_k (t) + Q_{rk} (x) U_k (t) + M_{rk} (x) U_k\,' (t) \biggr]


\equiv \sum \limits_{r=0}^{\infty} \e^r Y_{rk} (x,t), \ \; k=1,2,


\end{equation}


соответствующих корням $k_{1, 2} = \pm \,i\,\sqrt{x\wt{a} (x)}$


и


\begin{equation}


Y_3 (x, \e) \equiv Y_3 (x, t, \e) \equiv \sum_{r=0}^{\infty} \e^r Y_{r3} (x)


\equiv \sum \limits_{r=0}^{2} \e^r C_{r3} x^{\rho} \beta (x) +


\sum \limits_{r=3}^{\infty} \e^r Y_{r3} (x). 


\end{equation}





Проведем сужение в решениях (5.6) при $t = \Phi(x, \e)$ и учтем


равенства (3.2). Тогда первые два решения СВДУ (0.1) будут иметь


следующую структуру:


$$


Y_k (x, \Phi(x, \e), \e) = V_{k} (x, \e) \int


\limits_{\infty}^{\Phi(x, \e)} U_k (\tau) d \tau + Q_k (x, \e)


U_k (\Phi(x, \e)) + M_k (x, \e) \frac{d U_k (\Phi(x, \e))}{d


(\Phi(x, \e))},


$$


где множители при СОФ являются аналитическими функциями малого


параметра $\e >0$ и достаточно гладкие для всех $x \in [0,1]$.


Поскольку точка $x=0$ является стабильной точкой поворота, то


функции Эйри $U_k (t)$ являются ограниченными функциями для всех


$t\geq 0$.





Следовательно, СОФ содержащиеся в решениях СВДУ (0.1), являются


ограниченными функциями для всех $x \in \bigl[ 0,


\frac{\f(1)}{\e} \bigr]$, кроме СОФ $ W''_k (t) \equiv U'_k (t)$,


$k =1,2$.








\section{Оценка остаточных членов}





Запишем построенные решения (5.6) и (5.7) в виде


\begin{equation}


Y_k (x, t, \e) = Y_{p\,k}^* (x, t, \e) + \e^{p+1}


\xi_{(p+1)k} (x, t, \e), \quad k = \overline{1,3},


\end{equation}


где


\begin{equation}


Y_{p\,k}^* (x,t,\e)=\sum_{r=0}^p\, \e^r Y_{r\,k}(x, t), \quad k=1,2, \quad


Y_{p\,3}^* (x, \e) = \sum \limits_{r=0} ^p\, \e^r Y_{r\,3}(x) 


\end{equation}


--- частичные $p$-суммы, a $\e^{p+1} \xi_{(p+1)k} (x, t, \e)$


--- остаточные члены соответствующих рядов.


Легко проверить, что частичные суммы (6.2) удовлетворяют уравнениям


\begin{equation}


\bold{\wt{L}}_{\e} Y_{p\,k}^* (x, t, \e) = - \e^{p + 1} {\wt{\rho}}_{(p+1)k}


(x,t, \e), \, k = 1,2, \quad \bold{\wt{L}}_{\e} Y_{p\,3}^* (x, \e) =


h (x) - \e^{p + 1} {\wt{\rho}}_{(p+1)3} (x, \e). 


\end{equation}





Изучим структуру правых частей уравнений (6.3). С


учетом результатов действия операторов $ \bold{R}_k$ на элементы


ПБР (3.5) имеем


\begin{gather}


{\wt{\rho}}_{(p+1)k} (x,t, \e) \equiv A_{(p+1)k}\,^* (x, \e) W_k


(t) + B_{(p+1)k}\,^* (x, \e) U_k (t) +D_{(p+1)k}\,^* (x, \e) U_k\, '


(x,t, \e), \ \ k=1,2,\\


%


\wt{\rho}_{(p+1)3} (x,t, \e) \equiv \wt{\rho}_{(p+1)3} (x, \e) \equiv


\sum_{s=0} ^2\, \e^s \, {\omega}_{(p-2+s) }\, ''' (x), 


\end{gather}


где коэффициенты при СОФ в (6.4) и $\omega_{(p+1)}\,^*(x,\e)$


являются достаточно гладкими функциями при $x \in [0,1]$ и 


многочленами не выше второй степени относительно малого


параметра $\e > 0$.





Поскольку выражения (6.4) и (6.5) ограничены в ПБР $Y_r$, то


построенные решения (5.6) и (5.7) являются формальными


асимптотическими рядами решений расширенного уравнения (2.3), а их


сужения --- формальными асимптотическими рядами решений СВДУ (0.1)


([6], с.\,52; [14], с.\,10).





Подставив (6.1) в (2.3),


для определения функций $ \xi_{(p+1)k} (x,t, \e)$ получим


расширенные уравнения


\begin{equation}


\bold{\wt{L}}_{\e} \xi_{(p+1)k} (x,t, \e) =


{\wt{\rho}}_{(p+1)k} (x,t, \e), \quad k=\overline{1,3}. 


\end{equation}





Для того чтобы получить соответствующие оценки для остаточных


членов решений (6.1), применим следующую методику. Правые части


уравнений (6.6) принадлежат соответственно подпространствам


$Y_{rk}$, $k = 1,2$, и $X_r$, и эти пространства инвариантны


относительно расширенного оператора $\bold{\wt{L}}_{\e}$.


Поэтому решения уравнений (6.6) тоже ищем соответственно в этих


подпространствах, т.\,е. в виде


\begin{gather}


\xi_{(p+1)k} (x,t, \e) \equiv V_{(p+1)k}\,^* (x, \e) W_k


(t) + Q_{(p+1)k}\,^* (x, \e) U_k (t) +M_{(p+1)k}\,^* (x, \e) U_k\,'


(x,t, \e), \ \ k=1,2, \\


%


\xi_{(p+1)3} (x,t, \e) \equiv \xi_{(p+1)3} (x, \e) \equiv


\omega_{(p+1) }\,^* (x, \e).


\end{gather}





Описанная процедура определения решений уравнений (6.6) в виде


(6.7) и (6.8) является обобщением классического метода


неопределенных коэффициентов при решении дифференциальных


уравнений. Поэтому этот метод будем называть {\it методом


неопределенных функций}.


Подставим (6.7) и (6.8) в (6.6). Тогда для


определения неизвестных коэффициентов этих решений


получим серии обыкновенных


дифференциальных уравнений (см. (3.9)--(3.12))


\begin{gather}


\bigg[\bold{L}_0 + \e^3 \frac{d^3}{d x^3} \bigg] V_{(p+1)k}\, ^* (x, \e)=


A_{(p+1)k}\, ^* (x, \e),


\quad \bigg[\bold{L}_0 + \e^3 \frac{d^3}{d x^3} \bigg]


\omega_{(p+1)3}\, ^* (x, \e) = \alpha_{(p+1)} \,^* (x, \e), \\


%


\bold{D}_1\, Q_{(p+1)k} \,^* (x, \e) - \varepsilon [\bold{d} +


\varphi (x)\, \bold{m}]\, M_{(p+1)k} \,^* (x, \e) +


\varepsilon^3\,\frac{d Q_{(p+1)k} \,^* (x, \e)}{d x^3}=


F_{(p+1)kQ}\,^* (x, \e), \notag \\


%


\bold{D}_2\, M_{(p+1)k} \,^* (x, \e)+ \varepsilon^2 \bold{m}


\, Q_{(p+1)k} \,^* (x, \e) + \varepsilon^3\,\frac{d M_{(p+1)k}\,^*


(x, \e)}{d x^3}= F_{(p+1)kM}\,^* (x, \e),


\end{gather}


в которых правые части являются известными, достаточно гладкими функциями для


всех $ x\in[0,1]$ и аналитическими относительно малого параметра $ \e > 0$.





Пусть $ X_{r \varepsilon}$ --- пространства бесконечно дифференцируемых


функций по переменной $ x \in [0,1]$ и аналитически зависящих от малого


параметра $\varepsilon > 0$.





В уравнениях (6.9) дифференциальный оператор по форме совпадает с


дифференциальным оператором уравнения (0.1), т.\,е. для получения


оценки остаточного члена имеем стандартную ситуацию ([6],


с.\,63, формула (2.85)). Однако перед нами сейчас стоит задача,


отличающаяся от начальной постановки решения уравнения (0.1),


которая состоит в следующем. На этом этапе необходимо


определить только регулярные части решения уравнения (0.1), т.\,е.


найти решения этих уравнений в пространстве $ X_{r \varepsilon}$.


Исследуем уравнения (6.9), (6.10) в этом пространстве.





1. Область определения всех возмущенных и вырожденных операторов


этих уравнений совпадаeт с пространством $ X_{r \varepsilon}$.





2. Каждый вырожденный оператор имеет обратный оператор для всех


$x \in [0,1]$ в нужном нам подпространстве. Здесь имеется в виду


требование существования только частных решений некоторых


дифференциальных уравнений из (6.9), (6.10).





Следовательно, к решениям уравнений (6.9) и (6.10) применим


метод малого параметра Пуанкаре--Ляпунова. Поскольку правые части


уравнений (6.9) и (6.10) являются функциями порядка $ O(1) \,


\varepsilon^{-1/4}$ относительно малого параметра, то и все


решения этих уравнений тоже будут того же порядка для всех $ x \in


[0,1]$.





Возвращаясь к функциям (6.7), (6.8) и учитывая свойства СОФ $W_k


(t)$ и их производных, получим асимптотические неравенства


\begin{equation}


\| \xi_{(p+1)k} (x, t, \e) \| \leq A_{(p+1)}


\e^{\frac{-1}{4}}, \quad p\geq 1,\quad


\| \xi_{(p+1)3} (x, \e) \|


\leq A_{(p+1)},\quad k=1,2,\quad p \geq 0, 


\end{equation}


где постоянная $\mathrm {A}_{(p +1)} $ не зависит от $ x \in [0,1]$


и малого параметра $\e > 0$.





С целью исследования предела полученного решения при $ \e


\rightarrow +0$ необходимо более


конкретно исследовать полученные оценки при $ p = 0$. Поскольку


$M_{0k} (x) \equiv 0$, т.\,е. отсутствуют СОФ $U_k \,' (t)$, то при


$ p = 0$ будут иметь место улучшенные оценки


\begin{equation}


\| \xi_{1k} (x, t, \e) \| \leq \mathrm{A}_{1}, \quad k =1,2,


\quad \| \xi_{13} (x, \e) \| \leq \mathrm{A}_{1}. 


\end{equation}





Проведем сужение в построенных решениях и соответствующих оценках


при $ t = \Phi (x, \e)$. Тогда получим формулы


\begin{gather}


Y_k (x, \Phi(x, \e), \e) = Y_{pk}^* (x, \Phi(x, \e), \e) +


\e^{p+1} \xi_{(p+1)k} (x, \Phi(x, \e), \e), \quad k =1,2,


\tag{$6.1^k$} \\


%


Y_3 (x, \e) = Y_{p3}^* (x, \e) + \e^{p+1} \xi_{(p+1)3} (x, \e)


\tag{$6.1^3$}


\end{gather}


и оценки


\begin{gather}


\| \xi_{(p+1)k} (x, \Phi(x, \e), \e) \|


\leq \mathrm{A}_{p+1} \e^{\frac{-1}{4}},


\quad k = 1, 2, \quad \| \xi_{(p+1)3} (x, \e) \| \leq \mathrm{A}_{p +1},


\quad p \geq 1, \tag{$6.11^0$} \\


%


\| \xi_{1k} (x,\Phi (x, \e), \e) \| \leq \mathrm{A}_1, \quad k


=1,2,\quad


\| \xi_{13} (x, \e) \| \leq \mathrm{A}_1,


\quad p =0, \tag{$6.12^0$}


\end{gather}


где постоянные $\mathrm{A}_{p +1}$ и $\mathrm{A}_1$ не зависят от


$x \in [0;1]$ и малого параметра $\e >0$.





С учетом полученных оценок следует, что построенные решения


(5.6) и (5.7) являются асимптотическими рядами решений расширенного


уравнения (2.3), а их сужения --- асимптотические ряды решений СВДУ


(0.1).





\begin{note}


Нами проведена оценка остаточных членов


асимптотики решения для однородного СВДУ (0.1). Если уравнение (0.1)


неоднородное, то в решении $ y_3 (x, \e)$ (см. (5.7)) появляется


еще одно слагаемое, зависящее только от переменной $x$ и не


зависящее от произвольных постоянных. В этом случае оценка


полученного решения $ y_3 (x, \e)$ по форме остается прежней,


т.\,е. также имеют место оценки (6.11), (6.12) и $(6.12^0)$.


\end{note}





Одним из критериев правильности полученного результата является


исследование предела построенного решения при стремлении малого


параметра к нулю. Используя


свойства СОФ на бесконечности и оценки $(6.12^0)$, легко


установить, что на любом компакте отрезка $I$, не


содержащем точки $x = 0$, имеет место предельное равенство


\begin{equation}


\lim_{\e \to 0} Y(x, \Phi(x, \e), \e) = \lim_{\e \to 0} \biggl[


\sum \limits_{i=1}^2 Y_i (x, \Phi(x, \e), \e) + y_3 (x, \e) \biggr]


= y_{30} (x) = \o_0 (x), 


\end{equation}


где $\o_0 (x)$ --- решение вырожденного уравнения $\bold{L}_0


\o_0 (x) = h (x)$, $ x\in (0,1]$.





Сформулируем в виде теоремы полученные результаты.





\begin{thmnonum}


Пусть для СВДУ $(0.1){:}$ {\em a)}~имеют место условия $1^0$ и $2^0;$


{\em b)}~$-\frac{b(x)}{a(x)}=\frac{\rho}{x} + \aa(x)$, где


$\rho \in N $.


Тогда при достаточно малых значениях параметра $\e > 0$





$1)$ описанным методом можно построить три решения


расширенного уравнения $(2.3)$, представимые в виде асимптотических рядов


$(5.6)$, $(5.7);$





$2)$ сужение этих решений при $t=\Phi(x,\e)$, т.\,е. ряды $(6.1^k)$,


$(6.1^3)$ суть асимптотические ряды решений СВДУ $(0.1);$





$3)$ для остаточных членов асимптотических рядов решений $(6.1)$,


$k = \overline{1,3}$, имеют


место оценки $(6.11)$, $(6.12)$, $(6.11^0)$, $(6.12^0);$





$4)$ на любом компакте отрезка $I$, не содержащем точки $x=0$,


имеет место предельное равенство $(6.13)$.


\end{thmnonum}





\begin{note}


Если же функции $\wt{a}(x)$


и $b(x)$ имеют один и тот же знак при $x \in I$, то полученные


результаты в общем случае уже не имеют места.


\end{note}








\begin{thebibliography}{99}





\bibitem{1}


Langer~R. {\it The asymptotic solutions of ordinary linear


differential equations of the second order with special reference


to a turning point} // Trans. Amer. Math. Soc. -- 1949. -- V.\,67. --


P.\,461--490.


\bibitem{2}


Langer~R. {\it On the asymptotic forms of ordinary linear


differential equations of the third order in a region containing a


turning point} // Trans. Amer. Math. Soc. -- 1955. -- V.\,80.


-- P.\,93--123.


\bibitem{3}


Langer~R. {\it The solutions of a class of ordinary linear


differential equations of the third order in a region containing a


multiple turning point} // Duke Math. J. -- 1956. -- V.\,23.


-- P.\,93--110.


\bibitem{4}


Langer~R. {\it Formal solutions and a related equation for a class of


fourth order differential equations of a hydrodynamic type} //


Trans. Amer. Math. Soc. -- 1959. -- V.\,92. -- P.\,371--410.


\bibitem{5}


Дородницын~А.А. {\it Асимптотические законы распределения


собственных значений для некоторых особых видов дифференциальных


уравнений второго порядка} // УМН. -- 1952. -- Т.\,7.


\No\,6. -- C.\,3--96.


\bibitem{6}


Ломов~С.А. {\it Введение в общую теорию сингулярных возмущений}. --


М.: Наука, 1981. -- 400~c.


\bibitem{7}


Бобочко~В.Н. {\it Асимптотическое интегрирование системы


дифференциальных уравнений с точкой поворота } // Дифференц.


уравнения. -- 1991. -- Т.\,27. -- \No\,9. -- С.\,1505--1515.


\bibitem{8}


Бобочко~В.Н. {\it Уравнение типа Орра--Зоммерфельда с двумя точками поворота


} // Дифференц. уравнения. -- 1992. -- Т.\,28. -- \No\,10. -- C.\,1559--1570.


\bibitem{9}


Бобочко~В.Н. {\it Внутренняя точка поворота в теории сингулярных


возмущений}


 // Укр. матем. журн. -- 1996. -- Т.\,48. -- \No\,7. -- C.\,876--890.


\bibitem{10}


Бобочко~В.Н. {\it Системы дифференциальных уравнений с


сильной точкой поворота} // Укр. матем. журн.


-- 1997. -- T.\,49. -- \No\,11. -- C.\,1543--1547.


\bibitem{11}


Бобочко~В.Н. {\it Система дифференциальных уравнений с нестабильной


точкой поворота}. -- Ред. журн. ``Изв. вузов. Математика. -- 1997. --


\No\,3. -- C.\,73''. -- Казань, 1997. -- Деп. в ВИНИТИ, \No\,366-B97.


\bibitem{12}


Бобочко~В.Н. { \it Система дифференциальных уравнений с


точкой поворота в случае недиагонализируемого предельного


оператора} // Дифференц. уравнения. -- 1998. -- T.\,34.


-- \No\,10. -- С.\,1304--1312.


\bibitem{13}


Вазов~В. {\it Асимптотические разложения решений


дифференциальных уранений }. -- М.: Мир, 1968. -- 464~с.


\bibitem{14}


Ильин~А.М. {\it Согласование асимптотических разложений


решений краевых задач}. -- М.: Наука, 1989. -- 336~с.


\bibitem{15}


Вишик~М.И., Люстерник Л.А.


{\it Регулярное вырождение и пограничный слой для линейных


дифференциальных уравнений с малым параметром } // УMH.


-- 1957. -- Т.\,12. -- Bып.\,5. -- С.\,3--122.





\end{thebibliography}





\vskip 0.5cm





\postup{\it Кировоградский государственный}{\it Поступили}


\postup{\it педагогический университет}{\it первый вариант $05.01.1999$}


\postup{\it }{\it окончательный вариант $29.05.2001$}





\label{end}


\end{document}





